%中文latax模版 (UTF-8 版本, 使用 xeCJK)
\documentclass[11pt,onecolumn,twoside]{kzllyy}
%%%%%%%%%%%%%%%%%%%%%%%%%%%%%%%%%%%%%%%%%%%%%%%%%%%%%%%%%%%%%%%%
%  packages
%    这部分声明需要用到的包
%%%%%%%%%%%%%%%%%%%%%%%%%%%%%%%%%%%%%%%%%%%%%%%%%%%%%%%%%%%%%%%%
\usepackage{xeCJK}         % xeCJK 中文支持 (UTF-8)
\usepackage{fancyhdr}
\usepackage{amsmath,amsfonts,amssymb,graphicx}    % EPS 图片支持
\usepackage{subfigure}   % 使用子图形
\usepackage{indentfirst} % 中文段落首行缩进
\usepackage{bm}          % 公式中的粗体字符(用命令\boldsymbol）
\usepackage{multicol}    % 正文双栏
\usepackage{indentfirst} % 中文首段缩进
%\usepackage{picins}      % 图片嵌入段落宏包 比如照片 (TeX Live 2026 已不自带)
\usepackage{abstract}    % 2栏文档，一栏摘要及关键字宏包
\usepackage{amsthm}      % 使用定理
\usepackage{booktabs}    % 使用表格
\usepackage{float}       % [H] 浮动体固定位置
\usepackage{titlesec}
\setmainfont{TeX Gyre Termes}  % Times替代，xelatex兼容
\usepackage{wasysym}
\usepackage[sort]{cite}
%\def\citedash{$\sim$} % for creating cites like [3,4],[7~10]
%\renewcommand{\theequation}{\thesection.\arabic{equation}}(公式按节编号)
%%%%%%%%%%%%%%%%%%%%%%%%%%%%%%%%%%%%%%%%%%%%%%%%%%%%%%%%%%%%%%%%
%  xeCJK 字体设置
%    根据系统可用字体配置中文字体
%%%%%%%%%%%%%%%%%%%%%%%%%%%%%%%%%%%%%%%%%%%%%%%%%%%%%%%%%%%%%%%%
\setCJKmainfont[BoldFont={SimHei},ItalicFont={KaiTi}]{SimSun}  % 宋体为主字体
\setCJKsansfont{SimHei}    % 黑体为无衬线字体
\setCJKmonofont{FangSong}    % 仿宋为等宽字体
\setCJKfamilyfont{song}{SimSun}    % 宋体
\setCJKfamilyfont{hei}{SimHei}     % 黑体
\setCJKfamilyfont{kai}{KaiTi}      % 楷体
\setCJKfamilyfont{fsong}{FangSong} % 仿宋
%%%%%%%%%%%%%%%%%%%%%%%%%%%%%%%%%%%%%%%%%%%%%%%%%%%%%%%%%%%%%%%%
%  lengths
%    下面的命令重定义页面边距，使其符合中文刊物习惯。
%%%%%%%%%%%%%%%%%%%%%%%%%%%%%%%%%%%%%%%%%%%%%%%%%%%%%%%%%%%%%%%%
\usepackage{lettrine}
\setlength{\voffset}{-25pt} \setlength{\topmargin}{0pt}
\setlength{\headsep}{1mm} \setlength{\textheight}{245mm}
\setlength{\textwidth}{171mm} \setlength{\hoffset}{4.9pt}
\setlength{\evensidemargin}{-29pt}\setlength{\oddsidemargin}{0pt}
\setlength{\parskip}{1pt} \setlength{\columnsep}{7mm}
\setlength{\arraycolsep}{1pt}
\renewcommand{\baselinestretch}{1.15} %定义行间距
\setcounter{page}{1} \makeatletter
\renewcommand\section{\@startsection {section}{1}{\z@}%
                                   {-1pt }%
                                   {1pt}%
                                   {\normalfont\large\bfseries}}
\renewcommand\subsection{\@startsection {subsection}{1}{\z@}%
                                   {-1pt }%
                                   {1pt}%
                                   {\normalfont\normalsize\bfseries}}
\renewcommand\subsubsection{\@startsection {subsubsection}{1}{\z@}%
                                   {-1pt }%
                                   {1pt}%
                                   {\normalfont\normalsize\bfseries}}
%%%%%%%%%%%%%%%%%%%%%%%%%%%%%%%%%%%%%%%%%%%%%%%%%%%%%%%%%%%%%%%%
\setlength{\parskip}{1pt} \allowdisplaybreaks \makeatother
\newtheoremstyle{mystyle}{3pt}{3pt}{}{\parindent}{\bfseries}{}{5mm}{}
\theoremstyle{mystyle}
\newtheorem{Cor}{\normalsize{{\CJKfamily{hei}推论}}}
\newtheorem{Thm}{\normalsize{{\CJKfamily{hei}定理}}}
\newtheorem{Site}{\normalsize{{\CJKfamily{hei}引理}}}
\newtheorem{Def}{\normalsize{{\CJKfamily{hei}定义}}}
\newtheorem{Rem}{\normalsize{{\CJKfamily{hei}注}}}
\newtheorem{Exa}{\normalsize{{\CJKfamily{hei}例}}}
\newtheorem{Ste}{\normalsize{{\CJKfamily{hei}步骤}}}
\newtheorem{Sup}{\normalsize{{\CJKfamily{hei}假设}}}
\newtheoremstyle{citing}{3pt}{3pt}{}{}{\bfseries}{.}{.5em}{\thmnote{#3}}
\theoremstyle{citing}\newtheorem*{citedthm}{}
%%%%%%%%%%%%%%%%%%%%%%%%%%%%%%%%%%%%%%%%%%%%%%%%%%%%%%%%%%%%%%%%
\newcommand{\diag}{{\rm diag}}
\newcommand{\rank}{{\rm rank}}
\newcommand{\tr}{{\rm tr}}
\newcommand{\sign}{{\rm sign}}
\newcommand{\sgn}{{\rm sgn}}
\newcommand{\Proj}{{\rm Proj}}
\makeatletter
\renewcommand{\proofname}{\normalsize{{\CJKfamily{hei}证}}}
\renewenvironment{proof}[1][\proofname]{\par
  \pushQED{\qed}%
  \normalfont \topsep.1\p@\@plus.1\p@ \labelsep0em\relax
  \trivlist
  \item[\hskip 2.2em
        \bfseries
    #1]\hskip .5em\ignorespaces
}{%
  \endtrivlist\@endpefalse
} \makeatother

%%%%%%%%%%%%%%%%%%%%%%%%%%%%%%%%%%%%%%%%%%%%%%%%%%%%%%
%define the format like [1-2]
\makeatletter
\def\@compress@cite#1{%  % This is executed for each number
  \advance\@tempcnta\@ne % Now \@tempcnta is one more than the previous number
  \ifnum #1=\@tempcnta   % Number follows previous--hold on to it
        \def\@h@ld{\citedash \citeform{#1}}%
  \else   %  non-successor -- dump what's held and do this one
     \@h@ld \@citea \citeform{#1}%
     \let\@h@ld\@empty
  \fi \@tempcnta#1\let\@citea\citepunct
} \makeatother
%%%%%%%%%%%%%%%%%%%%%%%%%%%%%%%%%%%%%%%%%%%%%%%%%%%%%%

%%%%%%%%%%%%%%%%%%%%%%%%%%%%%%%%%%%%%%%%%%%%%%%%%%%%%%%%%%%%%%%%
% 首页页眉页脚定义
%%%%%%%%%%%%%%%%%%%%%%%%%%%%%%%%%%%%%%%%%%%%%%%%%%%%%%%%%%%%%%%%
\fancypagestyle{plain}{ \fancyhf{}\vspace{15pt}
\lhead{\footnotesize \quad 第~xx~卷 第~x~期\\
{\quad xxxx~年~x~月}}
\chead{\centering{{\CJKfamily{kai}\large 控~~制~~理~~论~~与~~应~~用}\\
{{\small Control Theory \& Applications}}}}
\rhead{\footnotesize {Vol. xx No. x~~~~~}\\
{Xxx. xxxx~~~~}} \lfoot{} \cfoot{} \rfoot{}
\renewcommand{\headrule}{%
\hrule height0.4pt width \headwidth \vskip1.0pt%
\hrule height0.4pt width \headwidth \vskip-2pt}}
%\renewcommand{\footrulewidth}{0.4pt}
%\setlength{\hoffset}{0.1pt}
%%%%%%%%%%%%%%%%%%%%%%%%%%%%%%%%%%%%%%%%%%%%%%%%%%%%%%%%%%%%%%%%
\pagestyle{fancy} \fancyhf{} \fancyhead[RE]{{\footnotesize
第~xx~卷~~~}} \fancyhead[CE]{{\footnotesize
控~~~制~~~理~~~论~~~与~~~应~~~用}}
\fancyhead[LE,RO]{~~~\footnotesize\thepage ~~~}
\fancyhead[LO]{{\footnotesize ~~~第~x~期}} \lfoot{} \cfoot{}
\rfoot{}
